% Documento principal. Formato APA 7 según la carta de la Gobernación del Huila
% (Radicado 2026SAL00020915, 2026-03-10): interlineado doble, márgenes de 2,54 cm,
% Times 12 pt, sangría de 1,27 cm. Lo gestiona el skill gestor-latex.
% Compilar: `make` (borrador) o `make final` (falla si quedan pendientes).

\PassOptionsToPackage{spanish,es-noshorthands}{babel}
\documentclass[man,12pt,babel,biblatex,nolmodern]{apa7}

% Tipografía: clon métrico de Times (la carta admite Times New Roman 12 pt)
\usepackage{newtxtext}
\usepackage{csquotes}
\usepackage{booktabs}
\usepackage{graphicx}
\usepackage{xcolor}

\AtBeginDocument{\setlength{\headheight}{16pt}}
\graphicspath{{figuras/}}
\addbibresource{bibliografia/referencias.bib}

% Idioma de las etiquetas de la clase
\addto\captionsspanish{%
  \renewcommand{\abstractname}{Resumen}%
  \renewcommand{\refname}{Referencias}%
}
\AtBeginDocument{\renewcommand{\keywordname}{Palabras clave}}

% biblatex-apa fija \& (inglés) entre los dos últimos autores; en español APA 7 usa «y»
% y no lleva coma antes de ella.
\DeclareDelimFormat[parencite,nptextcite]{finalnamedelim}{\addspace\bibstring{and}\space}
\DeclareDelimFormat[bib,biblist]{finalnamedelim}{%
  \ifthenelse{\value{listcount}>\maxprtauth}{}{\addspace\bibstring{and}\space}}

% Marca de pendiente: visible en modo borrador; error de compilación en modo final.
\newcommand{\pendiente}[1]{%
  \ifdefined\modofinal
    \PackageError{gestor-latex}{Pendiente sin resolver: #1}{Resuelva la marca o compile en modo borrador.}%
  \else
    \textbf{\textcolor{red}{[#1]}}%
  \fi}

\title{\pendiente{DATO PENDIENTE: título; por definir con el autor (P-018)}}
\shorttitle{Título pendiente}
\authorsnames{Jesús Santiago Ramón Ramos}
\authorsaffiliations{Programa de Ingeniería de Sistemas, Gobernación del Huila — Fondo Jenaro Díaz Jordán}
\authornote{La presente investigación se realiza como requisito para la condonación del crédito educativo otorgado por el Fondo Jenaro Díaz Jordán de la Gobernación del Huila, en concordancia con lo establecido en el Reglamento Operativo del año 2019. \pendiente{DATO PENDIENTE: alineación con el Plan de Desarrollo Departamental 2024-2027 (P-030)}}
\course{Neiva, Huila}
\date{2026}

\abstract{\pendiente{DATO PENDIENTE: resumen; se redacta al final con el skill redactar-resumen}}
\keywords{\pendiente{DATO PENDIENTE: palabras clave}}

\begin{document}

\maketitle

% GENERADO desde borrador_reescrito/01_introduccion.md — no editar a mano
\section{Introducción}

\subsection{Planteamiento del problema}

\pendiente{DATO PENDIENTE: planteamiento del problema; en discusión con el autor (P-028)}

\subsection{Pregunta de investigación}

\pendiente{DATO PENDIENTE: pregunta de investigación; por definir con el autor (P-018, P-029)}

\subsection{Justificación}

\pendiente{DATO PENDIENTE: justificación; se redacta cuando estén definidos el problema y la pregunta}

\subsection{Objetivos}

\subsubsection{Objetivo general}

\pendiente{DATO PENDIENTE: objetivo general; por definir con el autor (P-018)}

\subsubsection{Objetivos específicos}

\pendiente{DATO PENDIENTE: objetivos específicos; por definir con el autor (P-018)}

\subsection{Alcance y delimitación}

\pendiente{DATO PENDIENTE: alcance y límites del estudio; por definir con el autor (P-018)}

\subsection{Organización del documento}

\pendiente{DATO PENDIENTE: organización del documento; se redacta al cerrar la estructura}

% GENERADO desde borrador_reescrito/02_marco_teorico.md — no editar a mano

\section{Marco teórico}

\subsection{Marco conceptual}

Esta sección define los cuatro conceptos sobre los que se apoya el estudio: el trámite, la carga administrativa, la interoperabilidad y el principio de \enquote{solo una vez}. Se presentan en ese orden porque cada uno se apoya en el anterior: el trámite es el espacio en el que el ciudadano entrega información a la entidad; la carga administrativa describe el costo que esa entrega le supone; la interoperabilidad es el medio por el cual la entidad puede obtener la información sin pedirla, y el principio de \enquote{solo una vez} es el criterio que orienta ese cambio.

\subsubsection{Trámite}

En Colombia, el Decreto 088 de 2022 define el trámite como el \enquote{conjunto de requisitos, pasos o acciones, regulados por el Estado, dentro de un proceso misional, que deben efectuar los ciudadanos, usuarios o grupos de interés} ante una entidad pública o un particular que cumple funciones públicas, para hacer efectivo un derecho, ejercer una actividad o cumplir una obligación prevista o autorizada por la ley \parencite[art.~2.2.20.3, num.~13]{colombia2022decreto088}. Una definición cercana se encuentra en el estudio regional de \textcite{roseth2018complicada}, para quienes los trámites son el conjunto de requisitos, pasos o acciones por los cuales los individuos o las empresas \enquote{piden o entregan información a una entidad pública} con el fin de obtener un derecho o cumplir una obligación (p.~36); los autores indican que adaptaron esa definición de las de la Comisión Federal de Mejora Regulatoria de México y del \enquote{Departamento de Función Pública de Colombia} \parencite[p.~36, nota~1]{roseth2018complicada}. Ambas definiciones incluyen los requisitos como parte del trámite. A partir de ellas, en este estudio los documentos que la entidad exige al ciudadano se entienden como requisitos mediante los cuales este entrega información para obtener el derecho o cumplir la obligación.

Para analizar esos requisitos, el estudio los clasifica según quién expide o custodia el documento, con base en tres distinciones que hace la normativa colombiana. La primera agrupa los documentos que reposan en la propia entidad ante la cual se adelanta el trámite, que el Decreto Ley 019 de 2012 prohíbe exigir \parencite[art.~9]{colombia2012decreto019}. La segunda agrupa los documentos de competencia de otras autoridades, cuya presentación la Ley 962 de 2005 prohíbe solicitar \parencite[art.~1, num.~1]{colombia2005ley962}; dentro de ella se distinguen los que constan en registros públicos, que, mientras sus responsables no se integren al servicio ciudadano digital de interoperabilidad, deben poder consultarse en línea y de forma gratuita, de modo que el ciudadano queda eximido de aportar el certificado \parencite[art.~10, par.~3]{colombia2019decreto2106}. La tercera es una categoría residual, propuesta por este estudio, que reúne la información que solo el ciudadano puede aportar, como sus propias declaraciones o los formularios que diligencia, y que ninguna base de datos pública contiene. Además, \textcite{roseth2018complicada} señalan que los trámites \enquote{implican costos de transacción (algunos más que otros)} (p.~36) y reportan que, según los encuestados de Latinobarómetro en 2017, realizar su último trámite les tomó en promedio aproximadamente 5,4 horas en América Latina y el Caribe (p.~47). Ese costo es el objeto del concepto que se presenta a continuación.

\subsubsection{Carga administrativa}

Los costos que asume una persona al interactuar con el Estado se estudian bajo el concepto de carga administrativa. \textcite{moynihan2015administrative} la conceptualizan como una función de los costos de aprendizaje, psicológicos y de cumplimiento que experimentan los ciudadanos en sus interacciones con el Estado, y sostienen que el nivel de esa carga y su distribución entre el Estado y el individuo suelen ser resultado de decisiones políticas deliberadas y no simplemente de accidentes históricos o descuidos (p.~43). De los tres tipos de costo, los de cumplimiento son las cargas de seguir las reglas y requisitos administrativos; para quien solicita un servicio, son los costos de diligenciar formularios o de aportar documentos que acrediten su condición \parencite[p.~46]{moynihan2015administrative}. Trabajos posteriores de los mismos autores mantienen esa tipología y ubican el papeleo y la documentación dentro de los costos de cumplimiento \parencite{herd2025administrative} \pendiente{CONFIRMAR}. Este estudio no mide la carga que experimentan los ciudadanos del Huila; usa el concepto para interpretar cada documento que se exige como un costo de cumplimiento y, en consecuencia, dejar de exigirlo como una forma de trasladar esa parte de la carga a la administración. En ese sentido, la reducción de la carga administrativa es el beneficio esperado para el ciudadano que justifica la propuesta. El concepto siguiente describe el medio que hace posible ese traslado.

\subsubsection{Interoperabilidad}

El Marco de interoperabilidad para Gobierno Digital, publicado por el Ministerio de Tecnologías de la Información y las Comunicaciones (MinTIC), acoge como definición de interoperabilidad \enquote{la capacidad de las organizaciones para intercambiar información y conocimiento en el marco de sus procesos de negocio} \parencite[p.~33]{mintic2019marco}; según esa definición, el propósito del intercambio es que ciudadanos, empresas y otras entidades puedan recibir servicios en línea gracias a que los sistemas de las organizaciones comparten datos. El mismo documento describe que, en un escenario funcional de interoperabilidad, cuando una entidad requiera comprobar una circunstancia que obra en otra entidad, \enquote{no debe solicitarle dicha información al usuario, sino que deberá obtenerla de la entidad respectiva} por un mecanismo digital \parencite[p.~19]{mintic2019marco}. A ese mecanismo el marco lo denomina servicio de intercambio de información, y lo define como un recurso tecnológico que, mediante un conjunto de protocolos y estándares, permite el intercambio de información \parencite[p.~10]{mintic2019marco}. En la normativa colombiana, los Servicios Ciudadanos Digitales se definen como el conjunto de soluciones y procesos transversales que brindan al Estado capacidades para su transformación digital y para una adecuada interacción con el ciudadano \parencite[art.~2.2.17.1.4, num.~13]{colombia2020decreto620}, y entre los servicios base se incluye el servicio de interoperabilidad, que brinda las capacidades necesarias para garantizar el adecuado flujo de información entre los sistemas de las entidades \parencite[art.~2.2.17.2.1.1, num.~1.1]{colombia2020decreto620}. La descripción del marco corresponde a una orientación técnica y no constituye por sí misma una obligación legal; las obligaciones normativas se presentan en el marco normativo de este capítulo.

El marco de MinTIC señala, además, que la interoperabilidad no es un concepto meramente técnico, y la organiza en cuatro dominios: político-legal, organizacional, semántico y técnico \parencite[p.~33]{mintic2019marco}. La literatura académica describe ese mismo carácter amplio. \textcite{pardo2012interoperability} entienden la interoperabilidad en el gobierno electrónico como un conjunto de capacidades multidimensionales, complementarias y dinámicas que necesitan las redes de organizaciones para compartir información con éxito \pendiente{CONFIRMAR}. Ese carácter amplio también ayuda a explicar sus dificultades: \textcite{scholl2007integration} señalan que la mayoría de los esfuerzos de integración e interoperación enfrentan desafíos y restricciones importantes \pendiente{CONFIRMAR}. En Iberoamérica, \textcite{criado2011interoperabilidad} consultaron a los gobiernos de 16 países \parencite{criado2011interoperabilidad} y encontraron que el reto más mencionado para las iniciativas de interoperabilidad entre Estados era la cultura organizativa, señalada por el 87,5\,\% de los países consultados, y que, en el ámbito tecnológico, el principal problema eran los datos, en particular su dispersión en distintas bases y la ausencia de modelos de datos comunes \pendiente{DATO PENDIENTE: paginación de la revista para Criado et al. (2011)}. Aunque esos resultados corresponden a 2011 y a iniciativas entre países, en este estudio se toman como indicio de que los límites de la interoperabilidad pueden estar tanto en la organización como en los propios datos, y no solo en la tecnología. Por esa razón, los cuatro dominios del marco de MinTIC se usan en el estudio como criterio para clasificar cada documento exigido según lo que tendría que resolverse para dejar de pedirlo: una habilitación legal, un acuerdo entre entidades, un lenguaje común sobre los datos o una conexión técnica.

\subsubsection{Principio de \enquote{solo una vez}}

La idea de que el Estado no pida al ciudadano información que ya posee se conoce en la literatura como el principio de \enquote{solo una vez} (\emph{once-only principle}). \textcite{wimmer2021once} explica, en la introducción de su capítulo, que el principio figura entre los siete principios del plan de acción de administración electrónica 2016-2020 de la Unión Europea, con el fin de hacer el gobierno más eficaz y sencillo y de reducir la carga administrativa, pidiendo a ciudadanos y empresas que entreguen cierta información estándar a las autoridades una sola vez, y que el plan exige a las autoridades tomar medidas para reutilizar internamente esos datos, cuando esté permitido y con respeto de las normas de protección de datos. Su aplicación exige coordinar sistemas de varias entidades: el proyecto europeo con el que se buscó demostrarlo entre países se propuso conectar sistemas de información de varios países e identificar factores impulsores y barreras, como los requisitos de protección e intercambio de datos, los costos de implementación, el aislamiento entre entidades del sector público y los vacíos legales \parencite{krimmer2017exploring} \pendiente{CONFIRMAR}.

En América Latina y el Caribe, \textcite{roseth2018potencial} plantean que la interoperabilidad es lo que permite que el Estado comparta la información ciudadana con la que ya cuenta, de modo que el solicitante no la tenga que volver a presentar, y advierten que la región \enquote{en su mayoría está en una etapa incipiente en cuanto a la interoperabilidad} (p.~120). Según la encuesta BID-GEALC de 2017 que reportan los autores, 13 de 25 países de la región tenían iniciativas de \enquote{pedir una sola vez}, pero solo cuatro de esos 13 sabían cuántas entidades las cumplían \parencite[p.~147]{roseth2018potencial}. De esta manera, que el principio exista en una norma o en un programa no implica que se mida su cumplimiento.

En conjunto, los cuatro conceptos orientan los objetivos propuestos para el estudio, que siguen en revisión. La definición de trámite, la lectura de los documentos como requisitos y su clasificación según quién los expide sustentan la identificación de los documentos que exigen los trámites de la Gobernación del Huila; los conceptos de interoperabilidad, servicio de intercambio de información y Servicios Ciudadanos Digitales, junto con los cuatro dominios del marco de MinTIC, permiten relacionar esos documentos con las bases de datos y los mecanismos de intercambio disponibles y clasificar lo que haría falta para dejar de pedirlos; el principio de \enquote{solo una vez} orienta la propuesta, y la carga administrativa expresa el beneficio que se espera para el ciudadano. La sección siguiente presenta los marcos de interoperabilidad y las normas colombianas que permiten aplicar esos conceptos.

% GENERADO por gestor-latex (marcador) — sección sin redactar
\section{Metodología}

\pendiente{DATO PENDIENTE: sección sin redactar; usar el skill redactar-metodologia}

% GENERADO por gestor-latex (marcador) — sección sin redactar
\section{Resultados}

\pendiente{DATO PENDIENTE: sección sin redactar; usar el skill redactar-resultados}

% GENERADO por gestor-latex (marcador) — sección sin redactar
\section{Discusión}

\pendiente{DATO PENDIENTE: sección sin redactar; usar el skill redactar-discusion}

% GENERADO por gestor-latex (marcador) — sección sin redactar
\section{Conclusiones}

\pendiente{DATO PENDIENTE: sección sin redactar; usar el skill redactar-conclusiones}


\printbibliography[title={Referencias}]

\end{document}
